% The mathematical content of the blueprint, shared by the PDF and web builds.
% Each result carries:
%   \lean{...}   the Lean declaration that formalises it
%   \leanok      the Lean proof is complete and checked
%   \notready    stated here, Lean proof still to be written (Phase 2)
%   \uses{...}   the results it depends on (drives the dependency graph)

\chapter{The integer lattice base}

\begin{definition}[Sites]
  \label{def:site}
  \lean{DclFormalism.Site}
  \leanok
  For a dimension $d$, the \emph{sites} of the lattice $\Lambda_d$ are the
  integer points of $d$-dimensional space; that is, the elements of
  $\mathbb{Z}^d$.
\end{definition}

\begin{definition}[Parity colouring]
  \label{def:parity}
  \lean{DclFormalism.parity}
  \leanok
  \uses{def:site}
  The \emph{parity} of a site is the parity --- even or odd --- of the sum of
  its coordinates, taken as an element of the two-element ring
  $\mathbb{Z}/2\mathbb{Z}$. The even sites form one side of the bipartition,
  written $V^{+}$; the odd sites form the other, written $V^{-}$. This is the
  strict even/odd integer indexing of the diamond progression.
\end{definition}

\begin{definition}[Adjacency geometry]
  \label{def:geometry}
  \lean{DclFormalism.Geometry}
  \leanok
  \uses{def:site}
  A translation-invariant \emph{geometry} on $\Lambda_d$ is a finite set of
  neighbour offsets $O \subseteq \mathbb{Z}^d$. The \emph{neighbours} of a site
  $x$ are the translates $x + o$ for $o \in O$, and the \emph{coordination
  number} at $x$ is the number of distinct neighbours. Different geometries ---
  nearest-neighbour, extended-neighbour, the diamond progression --- are simply
  different offset sets, and so carry different coordination numbers. This is
  the general base that every specific lattice reuses.
\end{definition}

\begin{lemma}[Coordination is uniform]
  \label{lem:coord_card}
  \lean{DclFormalism.Geometry.coordination_eq_card}
  \leanok
  \uses{def:geometry}
  For every geometry $G$ and every site $x$, the coordination number at $x$
  equals the number of offsets $\lvert O \rvert$. In particular it is the same
  at every site.
\end{lemma}
\begin{proof}
  \leanok
  Translation by $x$ is injective, so the map $o \mapsto x + o$ sends the offset
  set bijectively onto the neighbour set and therefore preserves cardinality.
\end{proof}

\begin{definition}[Bipartite geometry]
  \label{def:bipartite}
  \lean{DclFormalism.Geometry.Bipartite}
  \leanok
  \uses{def:parity, def:geometry}
  A geometry is \emph{bipartite} (with respect to the parity colouring) when
  every offset has odd parity.
\end{definition}

\begin{lemma}[Edges cross the bipartition]
  \label{lem:neighbour_parity}
  \lean{DclFormalism.Geometry.neighbour_parity}
  \leanok
  \uses{def:bipartite}
  In a bipartite geometry every neighbour of a site has the opposite parity;
  hence every edge joins $V^{+}$ to $V^{-}$.
\end{lemma}
\begin{proof}
  \leanok
  Parity is additive under translation, so adding an odd-parity offset flips the
  coordinate-sum parity.
\end{proof}

\chapter{The diamond progression}

\begin{definition}[Diamond progression]
  \label{def:diamond}
  \lean{DclFormalism.IsDiamondProgression}
  \leanok
  \uses{def:bipartite, lem:coord_card}
  A geometry is a \emph{diamond progression} in dimension $d$ when it is
  bipartite and its coordination number equals $2d$ at every site.
\end{definition}

\begin{proposition}[Offset count]
  \label{prop:offsets_card}
  \lean{DclFormalism.IsDiamondProgression.offsets_card}
  \leanok
  \uses{def:diamond, lem:coord_card}
  A diamond-progression geometry has exactly $2d$ offsets.
\end{proposition}
\begin{proof}
  \leanok
  Immediate from the definition together with Lemma~\ref{lem:coord_card}: the
  coordination number, which the definition fixes at $2d$, equals the number of
  offsets.
\end{proof}

\chapter{The three result-density theorems}

These three theorems are the structural spine of Paper~IV. Their target values
are recorded in Lean (\texttt{coordinationTarget}, \texttt{maxHopTarget},
\texttt{unitCellVolumeTarget}); the proofs that the geometric quantities equal
them are the Phase-2 deliverable, so each statement is marked \notready\ until
its Lean proof lands.

\begin{theorem}[Coordination is linear]
  \label{thm:coord}
  \uses{def:diamond}
  \notready
  For the diamond progression, the coordination number is $\coord(d) = 2d$.
\end{theorem}

\begin{theorem}[Maximum hop]
  \label{thm:maxhop}
  \uses{def:diamond}
  \notready
  The longest single edge in the unit cell of the diamond progression has
  Euclidean length $\sqrt{d}$.
\end{theorem}

\begin{theorem}[Unit-cell volume]
  \label{thm:volume}
  \uses{def:diamond}
  \notready
  The unit cell of the diamond progression has volume $(d+1)^{(d-1)/2}$.
\end{theorem}
